% Generated draft: requires mathematical and source review.
\documentclass[11pt]{article}
\usepackage[margin=25mm]{geometry}
\usepackage{fontspec}
\setmainfont{Latin Modern Roman}
\usepackage{amsmath,amssymb,mathtools,unicode-math}
\setmathfont{Latin Modern Math}
\usepackage{enumitem,needspace,xurl,hyperref}
\setlength{\parindent}{0pt}
\setlength{\parskip}{5pt}
\setlength{\emergencystretch}{3em}
\newcommand{\N}{\mathbb{N}}
\newcommand{\Z}{\mathbb{Z}}
\newcommand{\Q}{\mathbb{Q}}
\newcommand{\R}{\mathbb{R}}
\newcommand{\C}{\mathbb{C}}
\newcommand{\F}{\mathbb{F}}
\newcommand{\A}{\mathbb{A}}
\newcommand{\PP}{\mathbb{P}}
\newcommand{\E}{\mathbb{E}}
\newcommand{\eps}{\varepsilon}
\newcommand{\dd}{\,\mathrm{d}}
\begin{document}
\section*{UnsolvedMath ID: 30005133}
\textbf{Source problem code:} catalogue entry
\subsection{Definitions and mathematical statement}
The supplied record is a theorem fragment, and no self-contained open problem can be recovered from it. The following displays are provisional transcriptions of its mathematical notation, not assertions of independently verified results:
\[
\pi_{*,0}\bigl(\mathbf{1}_{(p)}\bigr)(\F_p)\simeq\Z_{(p)},
\qquad
\operatorname{Map}_{\mathrm{SH}(\F_p)_{(p)}}(\mathbf{1},\mathbf{1})\simeq\Z_{(p)}.
\]
The source says that the second equivalence follows from the first and that an unspecified main theorem follows from the Hopkins--Morel conjecture. Neither that theorem nor the version of the conjecture used is supplied.

For orientation only, under conventional notation, if \(p\) is prime then \(\F_p=\Z/p\Z\) and
\[
\Z_{(p)}=\left\{\frac{a}{b}\in\Q: a,b\in\Z,\ b\ne0,\ p\nmid b\right\}.
\]
In motivic notation, \(\mathrm{SH}(k)\) usually denotes a stable motivic homotopy category over a field \(k\), and \(\mathbf{1}\) its monoidal unit. These conventions do not resolve the fragment: the record does not specify the quantifier on \(p\), the meaning of its parenthesized \((p)\) decorations, the grading of the right-hand side of the first display, or whether \(\operatorname{Map}\) denotes a mapping space, a mapping spectrum, or morphisms in a homotopy category. The placement of subscripts in these displays is itself an interpretation of the flattened source typography.

Consequently, no precise quantified conjecture is asserted here. The supplied verification describes the passage as explaining a proved main theorem, whereas the registry retains an \emph{open} label for an unclear extracted claim and marks its statement \emph{unrecoverable}. Recovering a mathematical target requires the surrounding report passage and its definitions.
\subsection{Short English statement}
The record contains two incompletely specified motivic homotopy equivalences and says that an unnamed main theorem follows from the Hopkins--Morel conjecture. It does not supply a self-contained question. Its retained \emph{open} label is a curation label; the supplied verification instead describes context for a proved theorem. The notation, hypotheses, and intended target require human review.
\subsection{Sources}
\begin{itemize}
\item Local UnsolvedMath registry: supplied record with permanent ID \(30005133\), exact source code \texttt{catalogue entry}, titled ``Hopkins--Morel Input to Motivic Stable Homotopy.'' No proposer is identified; the supplied proposed year is 2022.
\item Supplied report citation: \emph{Algebraic K-Theory (2023)}, Oberwolfach Reports, DOI 10.4171/owr/2022/24, pp.~1313--1376. Supplied source URL and literature-source link, labeled ``Oberwolfach report'': \url{https://doi.org/10.4171/owr/2022/24}. No precise page within this range or named author of the passage is supplied.
\item Literature assessment dated exactly \texttt{2026-08-22}: the supplied curation classifies the unclear extracted claim as open because it is a theorem fragment rather than an independently triageable statement. This is supplied dated curation, not an independently verified literature assessment.
\item The supplied statement-verification note says that the official passage explains how a proved main theorem follows from the Hopkins--Morel conjecture, that it is not a separate unresolved problem, and that no clean statement is promoted. This note is reported as supplied; the report has not been independently consulted for this draft.
\end{itemize}
\end{document}
